\Qhead{4.9}{Акустоэлектрический эффект и акустоэлектронные устройства}{Взаимодействие волн пространственного заряда с акустическим полем, акустоэлектрический эффект. Принципы работы акустоэлектронных устройств (усилители ультразвука, линии задержки, фильтры, конвольверы, запоминающие устройства).}
\Subfile{4.9.1}{sections/04_generation/q09/s01}

\recap{Пьезоэлектрическое поле связывает звук с носителями заряда. Дрейф способен изменить знак энергообмена. Медленное распространение звука реализует задержки, пространственное взвешивание — фильтры, нелинейное взаимодействие — свёртки.}
\board{Пьезоэлектрические связи; знак $\omega-\bm k\cdot\bm v_d$; два ВШП и задержка $L/v_s$.}
\oral{Откуда берётся энергия усиленного звука при акустоэлектрическом усилении?}
\bait{Можно оставить необходимость нелинейности конвольвера: линейная суперпозиция двух сигналов не создаёт их произведение.}
\sources{LMU~\cite{acoustoelectric}: взаимодействие ПАВ с носителями; Esslinger и соавт.~\cite{esslinger}: прямой эффект; Weng~\cite{weng}: университетская аннотация конкретной памяти. Линейный дрейфовый отклик, баланс импульса и формула свёртки выведены здесь. Полный текст Кайно и точная реализация курса не получены; совпадение всех коэффициентов с ними не заявляется.}
